% =========================================================
% preamble.tex
% Penghubung konfigurasi. Tidak memuat materi pembelajaran.
%
% Urutan pemuatan bersifat mengikat:
%   packages     -> package LaTeX
%   styles       -> palet warna dan tata letak
%   commands     -> custom command
%   environments -> environment OBE (memakai warna dari styles)
%   listings     -> konfigurasi source code (memakai warna dari styles)
%
% Makro \BookRoot ditetapkan oleh main.tex:
%   "."     -> buku dikompilasi utuh (main.tex)
%   "../.." -> satu bab dikompilasi mandiri dari chapters/NN/
% =========================================================

\input{\BookRoot/config/packages}
\input{\BookRoot/config/styles}
\input{\BookRoot/config/commands}
\input{\BookRoot/config/environments}
\input{\BookRoot/config/listings}

% ---------------------------------------------------------
% Karakter kutip ganda pada babel indonesian bukan shorthand,
% sehingga tidak diperlukan \shorthandoff. Pengaturan ini
% disediakan sebagai pengaman bila bahasa babel diubah kelak.
% ---------------------------------------------------------
\AtBeginDocument{%
    \iflanguage{indonesian}{}{\shorthandoff{"}}%
}

% ---------------------------------------------------------
% Direktori sumber daya (kode dan aset)
% ---------------------------------------------------------
\newcommand{\CodeDir}{\BookRoot/code/asm}
\newcommand{\AssetDir}{\BookRoot/assets}

% Nama berkas bibliografi TIDAK boleh diawali './' atau '/'.
% BibTeX memperlakukan nama berawalan path sebagai lokasi pasti
% dan tidak mencarinya melalui BIBINPUTS, sehingga
% \bibdata{./references} akan gagal ketika BibTeX dijalankan
% dari direktori lain (misalnya direktori output).
\ifSubfilesClassLoaded{%
    \newcommand{\BibFile}{../../references}%
}{%
    \newcommand{\BibFile}{references}%
}

\graphicspath{%
    {figures/}%
    {assets/images/}%
    {assets/diagrams/}%
    {../../assets/images/}%
    {../../assets/diagrams/}%
}

% ---------------------------------------------------------
% Resolusi \input relatif terhadap berkas yang sedang dibaca.
%
% Package subfiles memuat package import, yang memelihara
% direktori berkas aktif pada makro \import@path.
%
% Dengan menambahkan \import@path ke \input@path, perintah
%     \input{rangkuman}
% di dalam chapters/NN/chapter.tex akan dicari pada direktori
% bab tersebut, baik ketika buku dikompilasi utuh maupun ketika
% bab dikompilasi mandiri.
%
% Hal ini diperlukan karena \input biasa selalu relatif
% terhadap direktori kerja, bukan terhadap berkas pemanggil.
% ---------------------------------------------------------
\makeatletter
\def\input@path{{\import@path}}
\makeatother
